\scanpage{67}
\mspage{54}
% archive sheet number: 66

pour deux ouverts $U'\subset V'$ de $I'$, \deleted{dans le \unsure{dia}} (où on peut supposer que $V'=I'_{/i'}$), dans le diagramme obtenu de $(\alpha\alpha)$, $v'_{U',V'}$ soit dans $C$. Soit $U=u^{-1}(U')$, $V=u^{-1}(V')$, \deleted{de sorte que par transit\unsure{.}}, et considérons le diagramme similaire à $(\alpha\alpha)$, défini par les deux inclusions d'ouverts $U'\subset V'$, $U\subset V$. Cette fois, comme la formation de $u_!$ \unsure{commute} aux restrictions aux ouverts, \inserted{par transitivité des $\underline{\lim}$} les trois flèches verticales $c^X_{U'}$, $c^Y_{U'}$, $c^X_{V'}$ sont iso, donc aussi $c^f_{U',V'}$.

De $=$ la flèche $Y_V\mathbin{-}Y'_{V'}$ (à la place de $\unsure{Y|X}(i)\mathbin{-}Y'(i')$) est iso, Comme $v_{U,V}\in C$, il en est de $=$ de $v_{U',V'}$, ce qui prouve que $f'\in C_{I'^{\unsure{\circ|0}}}$. \deleted{Il \unsure{en sera}} On \unsure{procède} de $=$ pour $TC_I$. Enfin, l'assertion que $u_!$ applique $M(I^\circ)_C$ dans $M(I'^\circ)_{C\unsure{*}}$ résulte des \deleted{fait} définitions et de $u_!(C_{I^\circ})\subset C_{I'^\circ}$.

\uline{Corollaire} 3.3. (de Prop. 3.2. et Th. 3.1)

Soit $u:I\longrightarrow I'$ comme dans 3.2. Supposons de plus $I$ \deleted{\unsure{locm}} artinien, et
